\section{Conclusões}
\label{conclusoes}

% ORÇAMENTO: até 1½ coluna.
% Estrutura sugerida pelo template: (i) conclusão (hipóteses comprovadas? melhor método?),
% (ii) contribuições, (iii) trabalhos futuros.
% Retomar o objetivo geral e os específicos e dizer como foram (ou não) alcançados.
% Limitações principais (corredor indivisível superestima a distância, layout modelado, dados de outra realidade,
% baseline, escopo; layout por coocorrência ignora restrições reais).
% Trabalhos futuros: ex. posição do produto no corredor (retorno, maior lacuna), batching, restrições de precedência, Ratliff-Rosenthal, dados reais de uma loja.

% Papéis dos integrantes: as Diretrizes exigem que fique claro no relatório.
% O template não tem lugar para isso; confirmar com o professor se esta forma serve.
\section*{Contribuição dos Integrantes}
% Papel = os da Tabela 1 das Diretrizes. "Principal" e "Corresp." = atividades críticas em que
% o integrante foi responsável principal ou corresponsável. Manter curto (cabe em ~1/3 de coluna).

\begin{table}[!ht]
\caption{Papéis e responsabilidades dos integrantes no projeto.}
\label{tab-papeis}
\footnotesize
\begin{tabular}{@{}p{0.17\columnwidth}p{0.20\columnwidth}p{0.55\columnwidth}@{}}
\hline
\textbf{Integrante} & \textbf{Papel} & \textbf{Responsabilidades (principal / corresponsável)} \\ \hline
NomeA & Coordenador & ... \\
NomeB & Modelador A & ... \\
NomeC & Modelador B & ... \\
NomeD & Desenvolvedor A & ... \\
NomeE & Desenvolvedor B & ... \\ \hline
\end{tabular}
\end{table}
